\section{AHSME 1968}
        \begin{enumerate}
        	\item (\textit{AHSME 1968}) Let $P$ units be the increase in circumference of a circle resulting from an increase in $\pi$ units in the diameter. Then $P$ equals:


 $\text{(A) } \frac{1}{\pi}\quad\text{(B) } \pi\quad\text{(C) } \frac{\pi^2}{2}\quad\text{(D) } \pi^2\quad\text{(E) } 2\pi$


 

	\item (\textit{AHSME 1968}) The real value of $x$ such that $64^{x-1}$ divided by $4^{x-1}$ equals $256^{2x}$ is:


 $\text{(A) } -\frac{2}{3}\quad\text{(B) } -\frac{1}{3}\quad\text{(C) } 0\quad\text{(D) } \frac{1}{4}\quad\text{(E) } \frac{3}{8}$


 



 

	\item (\textit{AHSME 1968}) A straight line passing through the point $(0,4)$ is perpendicular to the line $x-3y-7=0$. Its equation is:


 $\text{(A) } y+3x-4=0\quad \text{(B) } y+3x+4=0\quad \text{(C) } y-3x-4=0\quad \\ \text{(D) } 3y+x-12=0\quad \text{(E) } 3y-x-12=0$


 

	\item (\textit{AHSME 1968}) Define an operation $\star$ for positive real numbers as $a\star b=\frac{ab}{a+b}$. Then $4 \star (4 \star 4)$ equals:


 $\text{(A) } \frac{3}{4}\quad \text{(B) } 1\quad \text{(C) } \frac{4}{3}\quad \text{(D) } 2\quad \text{(E )} \frac{16}{3}$


 

	\item (\textit{AHSME 1968}) If $f(n)=\frac{1}{3} n(n+1)(n+2)$, then $f(r)-f(r-1)$ equals:


 $\text{(A) } r(r+1)\quad \text{(B) } (r+1)(r+2)\quad \text{(C) } \frac{1}{3} r(r+1)\quad  \\ \text{(D) } \frac{1}{3} (r+1)(r+2)\quad \text{(E )} \frac{1}{3} r(r+1)(2r+1)$


 

	\item (\textit{AHSME 1968}) Let side $AD$ of convex quadrilateral $ABCD$ be extended through $D$, and let side $BC$ be extended through $C$, to meet in point $E.$ Let $S$ be the degree-sum of angles $CDE$ and $DCE$, and let $S'$ represent the degree-sum of angles $BAD$ and $ABC.$ If $r=S/S'$, then:


 $\text{(A) } r=1 \text{ sometimes, } r>1 \text{ sometimes}\quad\\ \text{(B) }  r=1 \text{ sometimes, } r<1 \text{ sometimes}\quad\\ \text{(C) } 0<r<1\quad \text{(D) } r>1\quad \text{(E) } r=1$


 

	\item (\textit{AHSME 1968}) Let $O$ be the intersection point of medians $AP$ and $CQ$ of triangle $ABC.$ if $OQ$ is 3 inches, then $OP$, in inches, is:


 $\text{(A) } 3\quad \text{(B) } \frac{9}{2}\quad \text{(C) } 6\quad \text{(D) } 9\quad \text{(E) } \text{undetermined}$


 

	\item (\textit{AHSME 1968}) A positive number is mistakenly divided by $6$ instead of being multiplied by $6.$ Based on the correct answer, the error thus committed, to the nearest percent, is :


 $\text{(A) } 100\quad \text{(B) } 97\quad \text{(C) } 83\quad \text{(D) } 17\quad \text{(E) } 3$


 

	\item (\textit{AHSME 1968}) The sum of the real values of $x$ satisfying the equality $|x+2|=2|x-2|$ is:


 $\text{(A) } \frac{1}{3}\quad \text{(B) } \frac{2}{3}\quad \text{(C) } 6\quad \text{(D) } 6\frac{1}{3}\quad \text{(E) } 6\frac{2}{3}$


 

	\item (\textit{AHSME 1968}) Assume that, for a certain school, it is true that


 I: Some students are not honest.
II: All fraternity members are honest.


 A necessary conclusion is:


 $\text{(A) Some students are fraternity members.} \quad\\ \text{(B) Some fraternity member are not students.} \quad\\ \text{(C) Some students are not fraternity members.} \quad\\ \text{(D) No fraternity member is a student.} \quad\\ \text{(E) No student is a fraternity member.}$


 

	\item (\textit{AHSME 1968}) If an arc of $60^{\circ}$ on circle $I$ has the same length as an arc of $45^{\circ}$ on circle $II$, the ratio of the area of circle $I$ to that of circle $II$ is:


 $\text{(A) } 16:9\quad \text{(B) } 9:16\quad \text{(C) } 4:3\quad \text{(D) } 3:4\quad \text{(E) } \text{none of these}$


 

	\item (\textit{AHSME 1968}) A circle passes through the vertices of a triangle with side-lengths $7\frac{1}{2},10,12\frac{1}{2}.$ The radius of the circle is:


 $\text{(A) } \frac{15}{4}\quad \text{(B) } 5\quad \text{(C) } \frac{25}{4}\quad \text{(D) } \frac{35}{4}\quad \text{(E) } \frac{15\sqrt{2}}{2}$


 

	\item (\textit{AHSME 1968}) If $m$ and $n$ are the roots of $x^2+mx+n=0 ,m \ne 0,n \ne 0$, then the sum of the roots is:


 $\text{(A) } -\frac{1}{2}\quad \text{(B) } -1\quad \text{(C) } \frac{1}{2}\quad \text{(D) } 1\quad \text{(E) } \text{undetermined}$


 

	\item (\textit{AHSME 1968}) If $x$ and $y$ are non-zero numbers such that $x=1+\frac{1}{y}$ and $y=1+\frac{1}{x}$, then $y$ equals


 $\text{(A) } x-1\quad \text{(B) } 1-x\quad \text{(C) } 1+x\quad \text{(D) } -x\quad \text{(E) } x$


 

	\item (\textit{AHSME 1968}) Let $P$ be the product of any three consecutive positive odd integers. The largest integer dividing all such $P$ is:


 $\text{(A) } 15\quad \text{(B) } 6\quad \text{(C) } 5\quad \text{(D) } 3\quad \text{(E) } 1$


 

	\item (\textit{AHSME 1968}) If $x$ is such that $\frac{1}{x}<2$ and $\frac{1}{x}>-3$, then:


 $\text{(A) } -\frac{1}{3}<x<\frac{1}{2}\quad \text{(B) } -\frac{1}{2}<x<3\quad \text{(C) } x>\frac{1}{2}\quad\\ \text{(D) } x>\frac{1}{2} \text{ or} -\frac{1}{3}<x<0\quad \text{(E) } x>\frac{1}{2} \text{ or } x<-\frac{1}{3}$


 

	\item (\textit{AHSME 1968}) Let $f(n)=\frac{x_1+x_2+\cdots +x_n}{n}$, where $n$ is a positive integer. If $x_k=(-1)^k, k=1,2,\cdots ,n$, the set of possible values of $f(n)$ is:


 $\text{(A) } \{0\}\quad \text{(B) } \{\frac{1}{n}\}\quad \text{(C) } \{0,-\frac{1}{n}\}\quad \text{(D) } \{0,\frac{1}{n}\}\quad \text{(E) } \{1,\frac{1}{n}\}$


 

	\item (\textit{AHSME 1968}) Side $AB$ of triangle $ABC$ has length 8 inches. Line $DEF$ is drawn parallel to $AB$ so that $D$ is on segment $AC$, and $E$ is on segment $BC$. Line $AE$ extended bisects angle $FEC$. If $DE$ has length $5$ inches, then the length of $CE$, in inches, is:


 $\text{(A) } \frac{51}{4}\quad \text{(B) } 13\quad \text{(C) } \frac{53}{4}\quad \text{(D) } \frac{40}{3}\quad \text{(E) } \frac{27}{2}$


 

	\item (\textit{AHSME 1968}) Let $n$ be the number of ways $10$ dollars can be changed into dimes and quarters, with at least one of each coin being used. Then $n$ equals:


 $\text{(A) } 40\quad \text{(B) } 38\quad \text{(C) } 21\quad \text{(D) } 20\quad \text{(E) } 19$


 

	\item (\textit{AHSME 1968}) The measures of the interior angles of a convex polygon of $n$ sides are in arithmetic progression. If the common difference is $5^{\circ}$ and the largest angle is $160^{\circ}$, then $n$ equals:


 $\text{(A) } 9\quad \text{(B) } 10\quad \text{(C) } 12\quad \text{(D) } 16\quad \text{(E) } 32$


 

	\item (\textit{AHSME 1968}) If $S=1!+2!+3!+\cdots +99!$, then the units' digit in the value of S is:


 $\text{(A) } 9\quad \text{(B) } 8\quad \text{(C) } 5\quad \text{(D) } 3\quad \text{(E) } 0$


 

	\item (\textit{AHSME 1968}) A segment of length $1$ is divided into four segments. Then there exists a quadrilateral with the four segments as sides if and only if each segment is:


 $\text{(A) equal to } \frac{1}{4}\quad\\ \text{(B) equal to or greater than } \frac{1}{8} \text{ and less than }\frac{1}{2}\quad\\ \text{(C) greater than } \frac{1}{8} \text{ and less than }\frac{1}{2}\quad\\ \text{(D) equal to or greater than } \frac{1}{8} \text{ and less than }\frac{1}{4}\quad\\ \text{(E) less than }\frac{1}{2}$


 

	\item (\textit{AHSME 1968}) If all the logarithms are real numbers, the equality
$log(x+3)+log(x-1)=log(x^2-2x-3)$
is satisfied for:


 $\text{(A) all real values of }x \quad\\ \text{(B) no real values of } x\quad\\ \text{(C) all real values of } x \text{ except } x=0\quad\\ \text{(D) no real values of } x \text{ except } x=0\quad\\ \text{(E) all real values of } x \text{ except } x=1$


 

	\item (\textit{AHSME 1968}) A painting $18$" X $24$" is to be placed into a wooden frame with the longer dimension vertical. The wood at the top and bottom is twice as wide as the wood on the sides. If the frame area equals that of the painting itself, the ratio of the smaller to the larger dimension of the framed painting is:


 $\text{(A) } 1:3\quad \text{(B) } 1:2\quad \text{(C) } 2:3\quad \text{(D) } 3:4\quad \text{(E) } 1:1$


 

	\item (\textit{AHSME 1968}) Ace runs with constant speed and Flash runs $x$ times as fast, $x>1$. Flash gives Ace a head start of $y$ yards, and, at a given signal, they start off in the same direction. Then the number of yards Flash must run to catch Ace is:


 $\text{(A) } xy\quad \text{(B) } \frac{y}{x+y}\quad \text{(C) } \frac{xy}{x-1}\quad \text{(D) } \frac{x+y}{x+1}\quad \text{(E) } \frac{x+y}{x-1}$


 

	\item (\textit{AHSME 1968}) Let $S=2+4+6+\cdots +2N$, where $N$ is the smallest positive integer such that $S>1,000,000$. Then the sum of the digits of $N$ is:


 $\text{(A) } 27\quad \text{(B) } 12\quad \text{(C) } 6\quad \text{(D) } 2\quad \text{(E) } 1$


 

	\item (\textit{AHSME 1968}) Let $S_n=1-2+3-4+\cdots +(-1)^{n-1}n$, where $n=1,2,\cdots$. Then $S_{17}+S_{33}+S_{50}$ equals:


 $\text{(A) } 0\quad \text{(B) } 1\quad \text{(C) } 2\quad \text{(D) } -1\quad \text{(E) } -2$


 

	\item (\textit{AHSME 1968}) If the arithmetic mean of $a$ and $b$ is double their geometric mean, with $a>b>0$, then a possible value for the ratio $a/b$, to the nearest integer, is:


 $\text{(A) } 5\quad \text{(B) } 8\quad \text{(C) } 11\quad \text{(D) } 14\quad \text{(E) none of these}$


 

	\item (\textit{AHSME 1968}) Given the three numbers $x,y=x^x,z=x^{x^x}$ with $.9<x<1.0$. Arranged in order of increasing magnitude, they are:


 $\text{(A) } x,z,y\quad \text{(B) } x,y,z\quad \text{(C) } y,x,z\quad \text{(D) } y,z,x\quad \text{(E) } z,x,y$


 

	\item (\textit{AHSME 1968}) Convex polygons $P_1$ and $P_2$ are drawn in the same plane with $n_1$ and $n_2$ sides, respectively, $n_1\le n_2$. If $P_1$ and $P_2$ do not have any line segment in common, then the maximum number of intersections of $P_1$ and $P_2$ is:


 $\text{(A) } 2n_1\quad \text{(B) } 2n_2\quad \text{(C) } n_1n_2\quad \text{(D) } n_1+n_2\quad \text{(E) } \text{none of these}$


 

	\item (\textit{AHSME 1968}) 
                    \begin{figure}[H]
                        \centering
                        \includegraphics[width=0.5\textwidth]{Problem/AHSME_I_1968_31_Problem_Figure_1.png}
                    \end{figure}
                    


 In this diagram, not drawn to scale, Figures $I$ and $III$ are equilateral triangular regions with respective areas of $32\sqrt{3}$ and $8\sqrt{3}$ square inches. Figure $II$ is a square region with area $32$ square inches. Let the length of segment $AD$ be decreased by $12\frac{1}{2}$ \% of itself, while the lengths of $AB$ and $CD$ remain unchanged. The percent decrease in the area of the square is:


 $\text{(A)}\ 12\frac{1}{2}\qquad\text{(B)}\ 25\qquad\text{(C)}\ 50\qquad\text{(D)}\ 75\qquad\text{(E)}\ 87\frac{1}{2}$


 

	\item (\textit{AHSME 1968}) $A$ and $B$ move uniformly along two straight paths intersecting at right angles in point $O$. When $A$ is at $O$, $B$ is $500$ yards short of $O$. In two minutes they are equidistant from $O$, and in $8$ minutes more they are again equidistant from $O$. Then the ratio of $A$'s speed to $B$'s speed is:


 $\text{(A) } 4:5\quad \text{(B) } 5:6\quad \text{(C) } 2:3\quad \text{(D) } 5:8\quad \text{(E) } 1:2$


 

	\item (\textit{AHSME 1968}) A number $N$ has three digits when expressed in base $7$. When $N$ is expressed in base $9$ the digits are reversed. Then the middle digit is:


 $\text{(A) } 0\quad \text{(B) } 1\quad \text{(C) } 3\quad \text{(D) } 4\quad \text{(E) } 5$


 

	\item (\textit{AHSME 1968}) With $400$ members voting the House of Representatives defeated a bill. A re-vote, with the same members voting, resulted in the passage of the bill by twice the margin by which it was originally defeated. The number voting for the bill on the revote was $\frac{12}{11}$ of the number voting against it originally. How many more members voted for the bill the second time than voted for it the first time?


 $\text{(A) } 75\quad \text{(B) } 60\quad \text{(C) } 50\quad \text{(D) } 45\quad \text{(E) } 20$


 

	\item (\textit{AHSME 1968}) 
                    \begin{figure}[H]
                        \centering
                        \includegraphics[width=0.5\textwidth]{Problem/AHSME_I_1968_35_Problem_Figure_1.png}
                    \end{figure}
                    


 In this diagram the center of the circle is $O$, the radius is $a$ inches, chord $EF$ is parallel to chord $CD$. $O$,$G$,$H$,$J$ are collinear, and $G$ is the midpoint of $CD$. Let $K$ (sq. in.) represent the area of trapezoid $CDFE$ and let $R$ (sq. in.) represent the area of rectangle $ELMF.$ Then, as $CD$ and $EF$ are translated upward so that $OG$ increases toward the value $a$, while $JH$ always equals $HG$, the ratio $K:R$ becomes arbitrarily close to:


 $\text{(A)} 0\quad\text{(B)} 1\quad\text{(C)} \sqrt{2}\quad\text{(D)} \frac{1}{\sqrt{2}}+\frac{1}{2}\quad\text{(E)} \frac{1}{\sqrt{2}}+1$


 

\end{enumerate}